\section*{Prologue}
\addcontentsline{toc}{section}{Prologue}

% TODO : vérifier l'ensemble des références bibliographiques avant la version finale.
% TODO : ajouter référence neurosciences cognitives (carte cognitive euclidienne / graphe cognitif)

% Ce texte part d'une question simple : qu'est-ce qu'une carte ?

% Avant de consulter un dictionnaire ou la littérature existante, j'ai essayé d'en construire moi-même une définition.

% La première intuition était qu'une carte est une représentation graphique. Mais cette propriété est manifestement insuffisante : une peinture est elle aussi une représentation graphique. Travaillant actuellement avec des cartographies bibliométriques, j'ai alors envisagé une seconde caractéristique : une carte représente des relations entre des éléments (des flux commerciaux, des vents, des frontières entre pays, etc). Cette formulation permet de distinguer la carte de nombreuses images, mais elle englobe également des objets que l'on ne qualifie généralement pas de cartes, comme les arbres généalogiques (c'est une représentation des relations génétiques au sein d'une famille).

% Une troisième hypothèse est alors apparue, issue de ma pratique. Lorsqu'une cartographie bibliométrique contient peu de n\oe{}uds et peu de connexions, j'aurais spontanément tendance à la décrire comme un arbre plutôt que comme une carte. La différence semblait donc pouvoir résider dans la densité des relations : peu de connexions produiraient un arbre, beaucoup de connexions une carte.

% Cette réponse paraissait plausible, mais elle résiste mal à l'examen. Elle suggère soit que les arbres constituent un cas particulier de carte, soit que la distinction entre les deux n'est qu'une question de degré. Surtout, elle fait reposer la frontière sur un critère imprécis dont l'évaluation dépend largement du contexte et du jugement de l'observateur.

% En cherchant des travaux susceptibles d'éclairer cette question, je suis rapidement tombé sur une littérature abondante consacrée aux dimensions politiques, rhétoriques ou historiques des cartes. Ces approches sont importantes, mais elles ne répondaient pas directement à l'interrogation structurelle qui m'intéressait. C'est à partir de ce constat que j'ai entrepris la réflexion présentée ici.

% Pour explorer la question, j'ai choisi de travailler en dialogue avec une intelligence artificielle. Ce choix répondait avant tout à une contrainte pratique : disposer rapidement de références, de synthèses bibliographiques et de pistes conceptuelles susceptibles d'orienter l'enquête. L'outil a servi d'interlocuteur de travail plutôt que d'autorité\footnote{Les systèmes d'intelligence artificielle ont également été utilisés lors de la rédaction de ce texte pour discuter certaines formulations, améliorer la clarté de passages particuliers et suggérer des pistes d'organisation. Les choix rédactionnels, les arguments développés et les éventuelles erreurs demeurent toutefois de la responsabilité de l'auteur.}. Les analyses proposées ont été discutées, reformulées ou écartées lorsqu'elles ne me paraissaient pas convaincantes. Les arguments développés dans ce texte n'engagent donc que leur auteur.

% La réflexion développée ici repose sur une intuition précise. La littérature en théorie cartographique s'est largement intéressée au pouvoir rhétorique et politique de la carte \parencite{harley2001, wood1992}, ainsi qu'à la relation entre l'objet cartographique et son lecteur \parencite{denil2021}. Les neurosciences cognitives distinguent quant à elles la carte cognitive euclidienne du graphe cognitif, deux formes de représentation spatiale associées à des mécanismes distincts. L'hypothèse explorée ici est qu'il existe un lien structurel entre ces trois niveaux : la métrique globale cohérente d'un espace est ce qui rend cet espace habitable, et c'est cette habitabilité qui ouvre l'interprétation libre. La carte n'est pas un objet, c'est une dyade : un créateur qui réduit intentionnellement un espace, un lecteur qui s'y projette librement.

% Les pages qui suivent cherchent à développer et à justifier cette hypothèse.

% \section*{Prologue}
% \addcontentsline{toc}{section}{Prologue}

% TODO : vérifier l'ensemble des références bibliographiques avant la version finale.

Ce texte part d'une question simple : qu'est-ce qu'une carte ?

La question m'est venue en travaillant sur des cartographies bibliométriques. Ces représentations sont couramment désignées comme des cartes, alors même qu'elles ne décrivent ni un territoire géographique ni un espace physique au sens habituel du terme. Cette observation a conduit à une interrogation plus générale : qu'est-ce qui fait qu'une représentation mérite d'être qualifiée de carte ?

En cherchant des travaux susceptibles d'éclairer cette question, je suis rapidement tombé sur une littérature abondante consacrée aux dimensions politiques, rhétoriques ou historiques des cartes. Ces approches sont importantes, mais elles ne répondaient pas directement à l'interrogation conceptuelle qui m'intéressait. C'est à partir de ce constat qu'a commencé la réflexion présentée ici.

Pour explorer cette question, j'ai choisi de travailler en dialogue avec une intelligence artificielle. Ce choix répondait avant tout à une contrainte pratique : disposer rapidement de références, de synthèses bibliographiques et de pistes conceptuelles susceptibles d'orienter l'enquête.\footnote{Les systèmes d'intelligence artificielle ont également été utilisés lors de la rédaction de ce texte pour discuter certaines formulations, améliorer la clarté de passages particuliers et suggérer des pistes d'organisation. Les choix rédactionnels, les arguments développés et les éventuelles erreurs demeurent toutefois de la responsabilité de l'auteur.} L'outil a servi d'interlocuteur de travail plutôt que d'autorité. Les analyses proposées ont été discutées, reformulées ou écartées lorsqu'elles ne paraissaient pas convaincantes. Les arguments développés dans ce texte n'engagent donc que leur auteur.

La réflexion développée ici repose sur une intuition précise. La littérature en théorie cartographique s'est largement intéressée au pouvoir rhétorique et politique de la carte \parencite{harley2001, wood1992}, ainsi qu'à la relation entre l'objet cartographique et son lecteur \parencite{denil2021}. Les neurosciences cognitives distinguent quant à elles la carte cognitive euclidienne du graphe cognitif, deux formes de représentation spatiale associées à des mécanismes distincts. L'hypothèse examinée dans les pages qui suivent est qu'il existe un lien structurel entre ces différents niveaux d'analyse : la cohérence métrique d'un espace pourrait être la condition qui le rend habitable pour un sujet et, par là même, interprétable. La carte n’est pas un objet, c’est une dyade : un créateur qui réduit intentionnellement un espace, un lecteur qui s’y projette librement.

Les pages qui suivent cherchent à développer et à justifier cette hypothèse.
